\textbf{Conciliación vigente INC-39.} La supervisión independiente añade cuatro EOLR-1 en extremos ordinarios, cuatro monitores y cuatro EOL a la selección desarrollada aquí: 56 monitores y 128 direcciones, cuatro SLC 40/40/24/24 con 88/88/52/52 mA. No cambia seis VEP, cuatro FRPS, ocho baterías 33 Ah ni cuatro 75 Ah. El controlador pasa a 1,232686/5,776 A y 36,078957 Ah de cribado; cada FRPS con relé a 0,625/0,751 A a 24 V y 18,0751 Ah. Los balances de INC-36 que siguen son la base antes de esos cuatro ramales, sin segunda lista activa. La selección y sus límites de subtensión, autonomía y supervisión están en INC/EXT-39.

\section{Distribución de energía de detección y actuación (INC-33)}
\label{sec:sd4-inc33}
INC-36 sustituye la distribución de INC-33 y es la selección vigente: seis VEP-A00-1P-UL activos, cuatro ordinarios y dos de bancos, tres repuestos fríos; cuatro FRPS activos, uno por VEP ordinario, y uno de reserva; ocho baterías 12 V/33 Ah activas y cuatro gabinetes, además de las reservas previas. Los dos VEP de bancos y seis DFIA quedan en Cheetah con SPS, respaldado por cuatro baterías 12 V/75 Ah en dos pares y dos gabinetes. El manual 06-356 rev.~11 establece 3 A en reposo y 12 A en alarma para el conjunto; los 2 A de salida auxiliar no amplían ese total \parencite[secs.~2.1.1, 3.4.2 y ap.~A]{lote33-cheetah11}.

El cálculo INC-33 de conexión directa anterior de cuatro VLF, dos VEP, seis DFIA, controlador y SPS alcanzaba 3,2474/9,035 A antes de módulos. Mostraba exceso de reposo y usaba 410 mA nominales VLF; no era máximo garantizado. Su distribución de dos FRPS con dos VLF y cribado de 24,7361 Ah por fuente queda reemplazada por un VEP por FRPS, con máximos 0,57/0,59 A más consumo propio y capacidad de cribado 17,4971 Ah. Las fichas VLF se mantienen como antecedente de esa sustitución, no como compra actual ni base del balance vigente \parencite[p.~2]{l18-vlf500ul}\parencite[p.~2]{lote30-vep36106}\parencite[ficha~06-432-1-5-DF]{lote30-fike542}\parencite[pp.~1--2]{lote33-frps}.

Los consumos VEP 0,57/0,59 A son máximos publicados a 24 V, no máximos demostrados en todo el intervalo de tensión. La salida AUX FRPS de 18,6--28,5 V es no supervisada y deja sólo 0,6 V frente al mínimo VEP de 18 V antes de caída; la supervisión del ramal y su tensión mínima conservan comprobación conforme a INC-36. Las cuatro FRPS seleccionadas admiten 1 A externo total en reposo y 10 A en alarma, con alimentación continua no rearmable al aspirante, falla individual supervisada y batería propia. Sus pares OEM de 33 Ah están dentro de los 35 Ah de carga admitidos; se comprueban curva útil, temperatura, envejecimiento, recarga y tensión final, sin declarar autonomía real por comparación nominal de Ah. Los cuatro circuitos AC a 240 V/50 Hz incorporan pérdidas y cargadores al balance pendiente; el subtotal de 58,08/70,176 W DC de las fuentes no es potencia de entrada AC ni eficiencia de catálogo \parencite[p.~2]{lote33-frps}\parencite[ap.~A, notas~3--6]{lote33-cheetah11}.

La lista INC-36 incluye 52 monitores 55-046, 54 relés 55-048, seis RCM y seis manuales más seis sensores presupuestados, dos de bancos aún por seleccionar: 124 direcciones. Sustituye el subtotal eléctrico de veinte monitores y la reserva de 122 direcciones. Se selecciona SLM 10-2473 y cuatro lazos con 40/40/22/22 direcciones y 88/88/48/48 mA de alarma frente a 100 mA por lazo. Seis avisadores P2RL-SP y su EOL/mapa NAC se desarrollan en INC-36. La corriente SLC presupuestada es 0,064346/0,272 A; el balance por fuente es controlador 1,230746/5,768 A y SPS 0,7954/3,45 A. Los cribados 36,0223/23,2525 Ah por fuente sustituyen el subtotal previo de 53,7826 Ah conjunto, sin acreditar mínima tensión al final de batería ni aptitud del confirmador de banco \parencite[ap.~A, secciones~1--3]{lote33-cheetah11}.

\section{Conexiones y compatibilidad de gas con descarga (EXT-33)}
\label{sec:sd4-ext33}
Se mantienen seis DFIA activos, uno por RCM 55-053; sus 24 V DC proceden de auxiliares protegidos del sistema, no del lazo SLC. La ficha DFIA permite conexión directa al RCM y exige el interruptor 02-12534 hacia módulo de monitorización Cheetah; no se inserta IRM. La señal IVOS identifica retirada física del actuador, mientras la presión acredita descarga: una no sustituye a la otra ni al estado de la válvula/compuerta. Cada VEP transmite fuego y falla por relés a circuitos supervisados independientes, conservando fallo por pérdida de alimentación y rotura de conductor; no se presume compatibilidad VEP/HLI por la sola existencia de un puerto VESDA \parencite[ficha~06-432-1-5-DF, pp.~5--6]{lote30-fike542}\parencite[secs.~3.8 y 4]{lote33-cheetah11}.

Para F04/F05 la condición de fuego mantiene alarma y retirada de recarga. Una falla de caudal, de medida H$_2$ o de su alimentación retira permiso sin esperar concentración alta. Antes de sellar se requieren clasificación del área en estados normal/falla/posaislamiento, emisor residual acotado y maniobra conjuntamente aprobada que conserve retención y alivio. El disparo no habilita recarga, ni el rearme del panel la repone. Mientras la condición de sellado no esté resuelta, no se presenta el diseño como extinción automática conforme de bancos. El conflicto exige cerrar diseño y cadena de seguridad antes de habilitar servicio, sin convertir la inhibición provisional en alternativa permanente a RT-06.17.

El cribado INT-31 de 810 m$^3$/h útiles, mezcla ideal y emisión cero reduce 0,8 a 0,2\,\% en $3600(36/810)\ln(4)=221,8071$ s; no es un temporizador de habilitación ni demuestra emisión nula tras retirar AC. Con recinto de 36 m$^3$ inicialmente a 0,2\,\%, diez minutos sellados y máximo 1,0\,\%, el límite propio sería $q_{res}\leq(0,010-0,002)(36)/(10/60)=1,728$ m$^3$/h; el límite conservador de intervención de 0,8\,\% daría 1,296 m$^3$/h. Ambos necesitan emisión residual realmente justificada y análisis de estratificación; no son datos del fabricante. Si se mantiene extracción de 810 m$^3$/h, un modelo ideal sin aporte de agente conserva tras diez minutos sólo $\exp[-810(10/60)/36]=0,0235177$ de la concentración inicial. Por ello mantener aquella extracción tampoco demuestra retención de diez minutos: se deben comprobar agente mínimo, fugas, alivio y evolución de gas conjuntamente. Los sectores ordinarios conservan su secuencia independiente, pendiente de hidráulica y selección completa.

\section{Placas, masa y segregación de pasos (PEN-33)}
\label{sec:sd4-pen33}
El detalle C-AJ-8309 requiere placa de al menos 6 mm en ambas caras del soporte; treinta aberturas independientes exigen sesenta placas y un dispositivo SF por cara, sesenta marcos SF6x1 SFF6000000121 y sesenta kits Wedge 120 AISI316 AWK0001201021. Los módulos y placas de relleno se dimensionan por cada cara; no se reutiliza un juego como sellado simultáneo de ambas. Se sustituyen las treinta unidades de PEN-27 por sesenta para este detalle, sin duplicar el número de cables que atraviesa cada vano. Para la reserva geométrica 379 por 476 mm, su masa propia de cribado es $60(0,379)(0,476)(0,006)(7850)=509,8217$ kg antes de marcos, módulos, anclajes y aberturas. No se duplica una piel que ya forme parte de placa OEM ni se compra una placa genérica bajo el listado. El valor PEN-30 de 254,9109 kg corresponde solamente a treinta placas: para dos caras independientes se usa el presente cálculo \parencite[pp.~1--4]{lote30-caj8309}.

Los doce paneles de 0,80 por 1,65 por 0,12 m aportan 4.561,92 kg brutos a 2.400 kg/m$^3$. Descontando treinta vanos de 201 por 298 por 120 mm, el neto geométrico de hormigón es 4.044,4013 kg; cada panel A con tres vanos resulta 328,4081 kg y cada B con dos 345,6588 kg. El presupuesto civil conserva la masa bruta como envolvente hasta armadura, encuentros y ejecución, sin descontar indebidamente barras ni sumarla otra vez al volumen neto. Con sesenta placas mínimas, hormigón neto más placas asciende a 4.554,2230 kg antes del resto. La armadura y anclajes se calculan por soporte y sismo; el contacto con BPR, aislamiento y juntas no hereda clasificación del paso. Se mantiene segregación A/B y pendiente lista de cables por material/diámetro: la denominación LSZH de CER-4 no equivale a las cubiertas PVC del listado, ni F2h acredita T2h de todos los servicios. El espesor de reserva con hormigón de 120 mm, dos placas de 6 mm y aislamiento de 350 mm es 482 mm dentro de los 500 mm previstos; excluye cualquier saliente del marco hacia la sala, que se comprueba en pasillos y soportes. Las longitudes de cables incorporan recorrido entre ambas caras, radios de curvatura y holguras de instalación sin fabricar empalmes dentro del paso. RT-06.04 conserva sus obligaciones de piso técnico/adaptación y enlaces dinámicos marinos.
